\bookchapter{Chapitre 11 : Photographie et gestion des images}

\chapterintro{Honnêtement, la plupart des gens prennent des photos au smartphone bien en deçà du potentiel de l'appareil. Ce n'est pas un problème de technique, mais de petits réglages et habitudes auxquels vous ne prêtez pas attention. Ce chapitre ne fait pas dans la théorie, juste ce que vous pouvez utiliser immédiatement.}

\booksection{Vous faites aussi ces erreurs ?}

Commençons par les problèmes les plus courants, voyez combien vous en commettez :

\begin{itemize}
  \item \textbf{L'objectif est sale}: C'est vraiment le problème le plus répandu. En sortant le smartphone, on met des empreintes sur l'objectif, les photos sont floues, on croit que le smartphone est en panne. Prenez l'habitude d'essuyer l'objectif avec un coin de vêtement avant de prendre une photo.
  \item \textbf{Ne jamais ajuster l'exposition}: En mode par défaut, taper l'écran fait la mise au point, mais si après avoir tapé vous glissez vers le haut ou le bas, vous pouvez régler la luminosité. Beaucoup de photos sont soit trop sombres soit trop claires simplement parce qu'on ne connaît pas cette manipulation.
  \item \textbf{Le flash toujours allumé}: Les photos au flash LED du smartphone sont vraiment laides, la lumière est dure et froide, on dirait des fantômes. Sauf s'il fait vraiment noir complet, désactivez-le. En faible luminosité, le mode nuit donne de bien meilleurs résultats que le flash.
  \item \textbf{Zoomer en écartant deux doigts}: Le pincement à deux doigts sur l'écran correspond au zoom numérique, qui recadre puis agrandit l'image, avec une baisse de qualité visible. Pour photographier de loin, utilisez directement les boutons de zoom du viseur (0,5x/1x/2x/3x). Parmi ceux-ci, 0,5x (ultra grand-angle) et 1x (objectif principal) sont des objectifs optiques natifs sans perte de qualité ; 2x/3x sont souvent des « zooms de qualité optique » — obtenus par recadrage du capteur haute résolution, la qualité est proche du zoom optique mais sans objectif dédié. Au-delà de 3x, il s'agit généralement de zoom hybride ou numérique, avec une baisse de qualité notable.
\end{itemize}

\booksection{Quelques astuces pour améliorer immédiatement vos photos}

Celles-ci ne nécessitent aucune connaissance théorique, un seul essai suffit pour voir le résultat :

\begin{itemize}
  \item \textbf{Règle des tiers + grille}: Activez la « grille » (aussi appelée lignes de référence) dans les paramètres de l'appareil photo. Placez ce que vous voulez mettre en valeur aux intersections des quatre lignes, l'image sera immédiatement « juste ». C'est la règle la plus simple et la plus efficace en photographie, sans ésotérisme, ça correspond simplement aux habitudes visuelles de l'œil humain.
  \item \textbf{Cherchez les « lignes directrices »}: En photographiant des paysages ou des rues, regardez s'il y a des lignes naturelles dans l'image — routes, rampes, murs, une rangée d'arbres — laissez-les s'étendre du bord de l'image vers le sujet. Avec cette technique bien maîtrisée, les passants penseront que vous avez étudié la photographie.
  \item \textbf{Ne photographiez pas toujours debout}: La plupart des gens prennent toujours les photos à la hauteur de leurs yeux, donc les perspectives sont toutes identiques. Essayez de vous accroupir pour un angle bas, ou de trouver un point en hauteur pour shooter vers le bas, le même endroit donne une sensation complètement différente. Pour les enfants et les animaux, il faut particulièrement s'accroupir, à hauteur des yeux pour un rendu naturel.
  \item \textbf{La lumière est le premier facteur}: À midi sous un soleil ardent, les ombres sur le visage sont dures et laides. La meilleure lumière est celle de « l'heure dorée » au petit matin et en fin de journée, le soleil est bas et la lumière chaude, tout est beau. La lumière d'un jour couvert n'est pas spectaculaire, mais elle est uniforme et douce, pour les portraits c'est souvent plus facile que sous un soleil direct.
  \item \textbf{Conseils pour le mode nuit}: Après activation du mode nuit, le smartphone prend automatiquement plusieurs photos et les fusionne. Le plus important alors — la main ne doit absolument pas bouger. Appuyez-vous contre un mur si possible, posez sur une table, ou à défaut retenez votre souffle. Le mode nuit exige que vous restiez stable 2 à 4 secondes, si vous bougez l'image sera floue.
\end{itemize}

\booksection{Mode portrait et mode professionnel}

\begin{itemize}
  \item \textbf{Mode portrait}: Il produit un effet d'arrière-plan flou avec un sujet net (faible profondeur de champ). Le smartphone ayant un objectif physiquement petit, les capacités réelles de flou d'arrière-plan sont limitées, il s'appuie donc sur des algorithmes IA — en résumé, l'IA vous détouré et floute le fond. À 1-2 mètres de distance l'effet est optimal, trop loin ou trop près l'algorithme a tendance à se tromper. De plus, en mode portrait, attention à ce que l'arrière-plan ne contienne pas de lignes trop complexes, sinon le détourage IA peut « déraper » sur les bords des cheveux.
  \item \textbf{Mode professionnel}: Si vous voulez un contrôle total sur le rendu de l'image, essayez le mode professionnel. Quelques paramètres clés à connaître : plus l'ISO est bas, plus l'image est propre, mais plus elle est sombre ; plus la vitesse d'obturation est lente, plus la lumière entre, mais le flou de bougé apparaît ; la balance des blancs gère le chaud/froid des couleurs, mal réglée la photo tire vers le bleu ou le jaune. La plupart des gens n'ont pas besoin du mode professionnel en permanence, mais connaître ces paramètres vous aide au moins à comprendre ce que fait l'appareil photo.
\end{itemize}

\booksection{Retouche photo : pas trop, juste ce qu'il faut}

Beaucoup de gens ont une idée fausse de la retouche, pensant que retoucher rend la photo artificielle. En réalité, la photo produite par l'appareil est déjà différente de ce que l'œil voit, un ajustement modéré sert à restituer ce que vous voyiez sur le moment, pas à tricher.

\begin{itemize}
  \item \textbf{Snapseed}: Outil de retouche gratuit de Google, fonctionnalités professionnelles sans publicité. Mon flux de retouche quotidien : d'abord « Ajuster l'image » pour la luminosité et le contraste « Détails » pour le sharpening « Balance des blancs » pour corriger les dominantes enfin « Recadrer » pour la composition voulue. La fonction « Local » de Snapseed est très pratique, elle permet de n'ajuster qu'une zone de l'image sans toucher au reste.
  \item \textbf{Xingtu/Meitu Xiuxiu}: Ces deux sont adaptés pour retoucher les portraits. Mais je pense qu'il y a un principe : le filtre beautifié en dessous de 30 \%, au-delà ça devient faux, la peau a l'air plastique. Le naturel est plus beau que la perfection.
  \item \textbf{Lightroom version mobile}: Si vous êtes exigeant sur les couleurs, c'est le plus professionnel. Il offre les réglages de courbes, le HSL par couleur et d'autres fonctions avancées. Mais pour la plupart des gens, c'est un peu utiliser un marteau-pilon pour écraser une mouche.
  \item \textbf{Conseil de flux de retouche}: Mon habitude est de d'abord prendre de bonnes photos originales, puis de les retoucher tranquillement à la maison sur l'écran de l'ordinateur ou de l'iPad, l'écran du smartphone est trop petit et les couleurs pas assez fiables. Si vous devez absolument retoucher sur smartphone, mettez la luminosité de l'écran à moyen-haut, ne retouchez pas en mode économie d'énergie.
\end{itemize}

\booksection{Trop de photos ? Stratégie d'organisation}

L'album du smartphone est rempli de milliers de photos, en chercher une précise prend une éternité — tout le monde connaît ça. Quelques méthodes pratiques :

\begin{itemize}
  \item \textbf{Utilisez la recherche}: Maintenant l'album de l'iPhone et de la plupart des smartphones Android prend en charge la recherche intelligente. Vous pouvez chercher « chat », « bord de mer », « carte d'identité », « fleurs » et le système utilise l'IA pour identifier le contenu de l'image et les retrouver. Vous pouvez aussi chercher par date et lieu, comme « mai dernier » ou « Pékin ». Peu de gens connaissent cette fonctionnalité, mais elle est vraiment très utile.
  \item \textbf{Nettoyage régulier des déchets}: Les captures d'écran, les photos en double, les ratés flous, ce sont eux les vrais responsables de l'encombrement. L'album iOS a une fonction de fusion automatique des « éléments en double » ; sur Android, allez dans le gestionnaire de fichiers et cherchez le dossier « screenshots » pour supprimer en masse les captures inutiles. Je passe cinq minutes par mois à nettoyer et je libère plusieurs Go.
  \item \textbf{Créez des albums, ne comptez pas uniquement sur le tri par défaut}: Le tri par défaut par date est très bien, mais pour les contenus qui traversent les périodes comme « les photos de croissance des enfants », « les photos de voyage », « les documents importants », vous devriez créer manuellement des albums pour les classer. L'« Album partagé » d'Apple permet aussi à toute la famille d'y ajouter des photos, une personne crée, toute la famille peut voir.
  \item \textbf{La sauvegarde est la ligne rouge}: Les photos sont les données les plus irremplaçables — les fichiers perdus peuvent être recréés, l'historique de discussion disparaît et c'est tout, mais les photos des enfants, les vidéos des grands-parents, une fois perdues c'est définitif. Activez absolument la sauvegarde cloud. iCloud Photos est le plus pratique, sur Android chaque marque a son propre service cloud, ou utilisez la sauvegarde automatique de Baidu Netdisk/Alibaba Cloud Drive. Le principe : au moins une sauvegarde cloud, ne mettez pas tous vos œufs dans le même panier.
\end{itemize}

\booksection{Capture d'écran et enregistrement d'écran}

\begin{itemize}
  \item \textbf{Capture d'écran}: Appuyez simultanément sur le bouton d'alimentation + volume haut, méthode universelle sur tous les smartphones. Beaucoup de smartphones Android prennent aussi en charge la capture par glissement à trois doigts vers le bas, ou le double tap avec l'articulation du doigt. La capture longue (capture défilante) est très mature sur Android — après une capture normale, cliquez sur le bouton « Défiler », le smartphone descend automatiquement les pages et les assemble en une seule image longue, idéale pour capturer des historiques de discussion et des pages web.
  \item \textbf{Enregistrement d'écran}: Ajoutez « Enregistrement d'écran » depuis le centre de contrôle, touchez pour commencer, touchez à nouveau pour arrêter. Avant d'enregistrer, appuyez longuement sur le bouton d'enregistrement pour choisir si vous enregistrez aussi le son du micro — si vous faites un tutoriel avec commentaires, activez le micro, si vous n'enregistrez que les manipulations à l'écran, désactivez-le. La vidéo enregistrée est automatiquement sauvegardée dans l'album.
\end{itemize}

\begin{tcbpractice}{Entraînez-vous : effectuez un cycle complet photo + retouche + organisation sur le smartphone}
\begin{enumerate}
  \item Activez la « grille » dans les paramètres de l'appareil photo, trouvez un endroit avec une lumière naturelle et prenez trois photos (n'oubliez pas d'essuyer l'objectif, ajustez l'exposition).
  \item Choisissez la meilleure, faites des ajustements de base avec Snapseed (luminosité, contraste, recadrage).
  \item Créez un album dédié pour cette photo, nommez-le « Mes œuvres ».
  \item Vérifiez que la sauvegarde cloud est activée, assurez-vous que la photo est déjà stockée dans le cloud.
\end{enumerate}

\end{tcbpractice}
